% R9_Structural_Classes.tex
% monogate Cost Theory — T41: Structural Classes Exhaustive, Mutually Exclusive,
% and Cost Ordering mean(C) < mean(B) < mean(A) < mean(D)
% Author: Arturo R. Almaguer
% Date: 2026-04-20
% Status: THEOREM (T41)

\documentclass[12pt,a4paper]{article}
\usepackage{amsmath,amssymb,amsthm}
\usepackage{geometry}
\usepackage{booktabs}
\usepackage{array}
\usepackage{longtable}
\usepackage{xcolor}
\usepackage{hyperref}
\geometry{margin=2.5cm}

% ── Theorem environments ────────────────────────────────────────────────────
\newtheorem{theorem}{Theorem}[section]
\newtheorem{lemma}[theorem]{Lemma}
\newtheorem{corollary}[theorem]{Corollary}
\newtheorem{proposition}[theorem]{Proposition}
\theoremstyle{definition}
\newtheorem{definition}[theorem]{Definition}
\newtheorem{conjecture}[theorem]{Conjecture}
\newtheorem{remark}[theorem]{Remark}
\newtheorem{example}[theorem]{Example}
\newtheorem{observation}[theorem]{Observation}

% ── Operator macros ─────────────────────────────────────────────────────────
\newcommand{\EML}{\operatorname{EML}}
\newcommand{\EDL}{\operatorname{EDL}}
\newcommand{\EXL}{\operatorname{EXL}}
\newcommand{\EAL}{\operatorname{EAL}}
\newcommand{\EMN}{\operatorname{EMN}}
\newcommand{\DEML}{\operatorname{DEML}}
\newcommand{\ELAd}{\operatorname{ELAd}}
\newcommand{\ELSb}{\operatorname{ELSb}}
\newcommand{\LEAd}{\operatorname{LEAd}}
\newcommand{\LEdiv}{\operatorname{LEdiv}}
\newcommand{\LEprod}{\operatorname{LEprod}}

% ── Cost macros ──────────────────────────────────────────────────────────────
\newcommand{\Cost}{\operatorname{Cost}}
\newcommand{\NC}{\operatorname{NaiveCost}}
\newcommand{\SD}{\operatorname{SharingDiscount}}
\newcommand{\PB}{\operatorname{PatternBonus}}
\newcommand{\Fam}{\mathcal{F}}
\newcommand{\cls}[1]{\textbf{Class~#1}}

% ── Cost shorthand ───────────────────────────────────────────────────────────
\newcommand{\cexp}{c_{\mathrm{exp}}}
\newcommand{\cln}{c_{\mathrm{ln}}}
\newcommand{\cmul}{c_{\mathrm{mul}}}
\newcommand{\cdiv}{c_{\mathrm{div}}}
\newcommand{\cadd}{c_{\mathrm{add}}}
\newcommand{\csub}{c_{\mathrm{sub}}}
\newcommand{\cpow}{c_{\mathrm{pow}}}
\newcommand{\cneg}{c_{\mathrm{neg}}}
\newcommand{\crec}{c_{\mathrm{rec}}}

% ── Column types ─────────────────────────────────────────────────────────────
\newcolumntype{L}[1]{>{\raggedright\arraybackslash}p{#1}}
\newcolumntype{C}[1]{>{\centering\arraybackslash}p{#1}}

\title{Theorem T41: Four Structural Classes Are Exhaustive,\\
       Mutually Exclusive, and Satisfy the Cost Ordering\\[6pt]
       \large $\overline{\Cost}_C < \overline{\Cost}_B
              < \overline{\Cost}_A < \overline{\Cost}_D$}
\author{Arturo R.\ Almaguer\\
\small monogate research \quad \texttt{almaguer1986@gmail.com}}
\date{April 2026}

% ════════════════════════════════════════════════════════════════════════════
\begin{document}
\maketitle

% ── Abstract ────────────────────────────────────────────────────────────────
\begin{abstract}
We prove Theorem~T41 of the monogate SuperBEST cost theory: the four
structural classes~A (Pure Exponential), B~(Rational), C~(Log-ratio),
and D~(Mixed) form an \emph{exhaustive} and \emph{mutually exclusive}
partition of the set of all finite arithmetic expressions computable
by EML trees over $\mathcal{F}_{16}$.
The partition is constructed by a two-bit signature
$(\mathbf{1}_{\mathrm{exp}},\mathbf{1}_{\mathrm{ln}})$ together with a
structural refinement for Class~D.
We then establish the cost ordering
$\overline{\Cost}_C < \overline{\Cost}_B < \overline{\Cost}_A
 < \overline{\Cost}_D$
by two complementary methods:
(i) a structural argument grounded in SuperBEST~v3 unit costs; and
(ii) an empirical confirmation across the 157-equation extended catalog,
reporting mean costs $\overline{\Cost}_C \approx 9.5$,
$\overline{\Cost}_B \approx 12.2$, $\overline{\Cost}_A \approx 10.4$,
$\overline{\Cost}_D \approx 20.1$ (50-equation COST-4 subset) and
consistent statistics from the enlarged catalog.
Ten worked classifications from the catalog illustrate the procedure.
Edge cases and borderline expressions are resolved in
Section~\ref{sec:edge}.
\end{abstract}

\tableofcontents
\bigskip

% ════════════════════════════════════════════════════════════════════════════
\section{Four Structural Classes: Formal Definitions}
\label{sec:classes}
% ════════════════════════════════════════════════════════════════════════════

Throughout, $E$ denotes a finite arithmetic expression computable by a
rooted EML tree over the 16-operator family $\mathcal{F}_{16}$.
The SuperBEST~v3 unit costs are fixed:
$\cexp = \cln = \cdiv = 1$;
$\cneg = \crec = \cmul = \csub = 2$;
$\cpow = \cadd = 3$.

% ── 1.1 Signature ──────────────────────────────────────────────────────────
\subsection{The Transcendental Signature}

\begin{definition}[Transcendental signature]
\label{def:sig}
For any expression~$E$ let
\[
  \sigma(E) \;=\;
  \bigl(\mathbf{1}[n_{\exp}(E) > 0],\;
        \mathbf{1}[n_{\ln}(E) > 0]\bigr)
  \;\in\; \{0,1\}^2,
\]
where $n_{\exp}(E)$ and $n_{\ln}(E)$ are the counts of
\texttt{exp} and \texttt{ln} nodes in the minimal EML tree for~$E$.
The four possible signatures are
$(0,0)$, $(1,0)$, $(0,1)$, and $(1,1)$.
\end{definition}

% ── 1.2 Class definitions ───────────────────────────────────────────────────
\subsection{Class Definitions}

\begin{definition}[Class B --- Rational]
\label{def:classB}
$E \in \mathbf{B}$ if and only if $\sigma(E) = (0,0)$: neither
\texttt{exp} nor \texttt{ln} appears.
$E$ is a rational function of its free terminals, possibly including
integer or rational powers.
\end{definition}

\begin{definition}[Class A --- Pure Exponential]
\label{def:classA}
$E \in \mathbf{A}$ if and only if $\sigma(E) = (1,0)$ (has \texttt{exp},
no \texttt{ln}) and every \texttt{exp} output in the parse tree of~$E$
satisfies both of the following conditions:
\begin{enumerate}
  \item Its argument is a rational sub-expression (no nested
        transcendentals in the argument).
  \item It is not combined via non-trivial arithmetic (\texttt{add},
        \texttt{sub}, or as numerator/denominator of a non-constant
        rational expression) with the output of any other \texttt{exp}
        node.
\end{enumerate}
Expressions with $\sigma(E) = (1,0)$ that violate either condition fall
into Class~D.
\end{definition}

\begin{definition}[Class C --- Log-ratio]
\label{def:classC}
$E \in \mathbf{C}$ if and only if $\sigma(E) = (0,1)$ (has \texttt{ln},
no \texttt{exp}) and the argument to every \texttt{ln} node is a rational
sub-expression (no transcendental nesting inside any \texttt{ln}
argument).
Expressions with $\sigma(E) = (0,1)$ that violate this condition fall
into Class~D.
\end{definition}

\begin{definition}[Class D --- Mixed]
\label{def:classD}
$E \in \mathbf{D}$ if and only if $E$ does not belong to~A, B, or~C.
Equivalently, $E \in \mathbf{D}$ whenever any of the following holds:
\begin{enumerate}
  \item $\sigma(E) = (1,1)$: both \texttt{exp} and \texttt{ln} present.
  \item $\sigma(E) = (1,0)$ but two or more \texttt{exp} outputs are
        combined at the same tree level via \texttt{add}, \texttt{sub},
        \texttt{div}, or \texttt{recip} with non-constant operands.
  \item $\sigma(E) = (0,1)$ but some \texttt{ln} argument contains a
        transcendental sub-expression.
\end{enumerate}
\end{definition}

\begin{remark}[Class D as complementary class]
Definitions~\ref{def:classA}--\ref{def:classC} give positive
characterisations of A, B, C.
Definition~\ref{def:classD} is the formal complement: $\mathbf{D}
= \mathcal{E} \setminus (\mathbf{A} \cup \mathbf{B} \cup \mathbf{C})$,
where $\mathcal{E}$ is the universe of finite EML-computable expressions.
\end{remark}

% ── 1.3 Cost ranges ─────────────────────────────────────────────────────────
\subsection{Characteristic Cost Ranges}

\begin{table}[h]
\centering
\caption{Structural classes: signature, cost range (50-equation COST-4
         catalog), mean cost, and representative examples.}
\label{tab:classes}
\smallskip
\begin{tabular}{@{}llcccl@{}}
\toprule
\textbf{Class} & \textbf{Signature}
  & \textbf{$n$}
  & \textbf{Cost range}
  & \textbf{Mean cost}
  & \textbf{Example} \\
\midrule
B: Rational          & $(0,0)$ & 19 & 2--34  & 12.2 & Michaelis--Menten \\
C: Log-ratio         & $(0,1)$ & 13 & 3--17  & 9.5  & pH formula \\
A: Pure Exponential  & $(1,0)$ & 8  & 5--13  & 10.4 & Arrhenius \\
D: Mixed             & ---     & 10 & 12--31 & 20.1 & Butler--Volmer \\
\bottomrule
\end{tabular}
\end{table}

% ════════════════════════════════════════════════════════════════════════════
\section{Exhaustive and Mutually Exclusive Partition}
\label{sec:partition}
% ════════════════════════════════════════════════════════════════════════════

\begin{theorem}[T41-Part: Classes A, B, C, D partition $\mathcal{E}$]
\label{thm:partition}
The four classes $\mathbf{A}$, $\mathbf{B}$, $\mathbf{C}$, $\mathbf{D}$
form a partition of the universe~$\mathcal{E}$ of finite
EML-computable arithmetic expressions:
\[
  \mathcal{E} = \mathbf{A} \sqcup \mathbf{B} \sqcup \mathbf{C}
              \sqcup \mathbf{D},
\]
where~$\sqcup$ denotes disjoint union.
Equivalently, every expression belongs to exactly one class.
\end{theorem}

The proof proceeds in two independent parts: exhaustiveness and mutual
exclusion.

% ── 2.1 Exhaustiveness ─────────────────────────────────────────────────────
\subsection{Exhaustiveness}

\begin{proposition}[Every expression is classified]
\label{prop:exhaust}
For every $E \in \mathcal{E}$, at least one of
$E \in \mathbf{A}$, $E \in \mathbf{B}$, $E \in \mathbf{C}$,
$E \in \mathbf{D}$ holds.
\end{proposition}

\begin{proof}
We case-split on $\sigma(E) \in \{0,1\}^2$.

\medskip
\noindent\textbf{Case 1: $\sigma(E) = (0,0)$.}
No \texttt{exp} and no \texttt{ln} in~$E$.
By Definition~\ref{def:classB}, $E \in \mathbf{B}$.

\medskip
\noindent\textbf{Case 2: $\sigma(E) = (1,1)$.}
Both \texttt{exp} and \texttt{ln} present.
By condition~(i) of Definition~\ref{def:classD}, $E \in \mathbf{D}$.

\medskip
\noindent\textbf{Case 3: $\sigma(E) = (1,0)$.}
\texttt{exp} present, \texttt{ln} absent.
We examine whether the two structural conditions of
Definition~\ref{def:classA} hold.
\begin{itemize}
  \item If both conditions hold: $E \in \mathbf{A}$.
  \item If condition~(1) fails (some \texttt{exp} argument contains
        a nested transcendental): $E \in \mathbf{D}$
        by condition~(ii) of Definition~\ref{def:classD}, since an
        \texttt{exp}-inside-\texttt{exp} structure represents two
        \texttt{exp} nodes in a non-trivial composition.
  \item If condition~(2) fails (two \texttt{exp} outputs combined by
        non-trivial arithmetic): $E \in \mathbf{D}$ by
        condition~(ii).
\end{itemize}
In all sub-cases, $E \in \mathbf{A} \cup \mathbf{D}$.

\medskip
\noindent\textbf{Case 4: $\sigma(E) = (0,1)$.}
\texttt{ln} present, \texttt{exp} absent.
We check whether every \texttt{ln} argument is rational.
\begin{itemize}
  \item If yes: $E \in \mathbf{C}$ by Definition~\ref{def:classC}.
  \item If no (some \texttt{ln} argument contains a transcendental
        nested within it): $E \in \mathbf{D}$ by condition~(iii)
        of Definition~\ref{def:classD}.
\end{itemize}
In all sub-cases, $E \in \mathbf{C} \cup \mathbf{D}$.

\medskip
\noindent Since $\{0,1\}^2 = \{(0,0),(1,1),(1,0),(0,1)\}$ is exhausted
and every case yields $E$ in at least one of $\mathbf{A}$, $\mathbf{B}$,
$\mathbf{C}$, $\mathbf{D}$, the union is exhaustive.
\end{proof}

% ── 2.2 Mutual exclusion ───────────────────────────────────────────────────
\subsection{Mutual Exclusion}

\begin{proposition}[No expression belongs to two classes]
\label{prop:exclusive}
The classes $\mathbf{A}$, $\mathbf{B}$, $\mathbf{C}$, $\mathbf{D}$
are pairwise disjoint.
\end{proposition}

\begin{proof}
We verify all six pairs.

\medskip
\noindent\textbf{$\mathbf{A} \cap \mathbf{B} = \emptyset$.}
$\mathbf{A}$ requires $n_{\exp} > 0$; $\mathbf{B}$ requires
$n_{\exp} = 0$.  Contradiction.

\medskip
\noindent\textbf{$\mathbf{A} \cap \mathbf{C} = \emptyset$.}
$\mathbf{A}$ requires $n_{\ln} = 0$; $\mathbf{C}$ requires
$n_{\ln} > 0$.  Contradiction.

\medskip
\noindent\textbf{$\mathbf{B} \cap \mathbf{C} = \emptyset$.}
$\mathbf{B}$ requires $n_{\ln} = 0$; $\mathbf{C}$ requires
$n_{\ln} > 0$.  Contradiction.

\medskip
\noindent\textbf{$\mathbf{A} \cap \mathbf{D} = \emptyset$.}
$\mathbf{A}$ requires both structural conditions of
Definition~\ref{def:classA}; $\mathbf{D}$ is entered (for the $(1,0)$
case) precisely when at least one of those conditions fails.
An expression cannot simultaneously satisfy and violate the same condition.

\medskip
\noindent\textbf{$\mathbf{B} \cap \mathbf{D} = \emptyset$.}
$\mathbf{B}$ requires $\sigma(E) = (0,0)$; no condition of
Definition~\ref{def:classD} is satisfiable when both
$n_{\exp} = 0$ and $n_{\ln} = 0$.
(Condition~(i) requires $n_{\ln} > 0$ and $n_{\exp} > 0$; conditions~(ii)
and~(iii) require at least one of $n_{\exp}$, $n_{\ln}$ to be positive.)

\medskip
\noindent\textbf{$\mathbf{C} \cap \mathbf{D} = \emptyset$.}
$\mathbf{C}$ requires all \texttt{ln} arguments to be rational;
$\mathbf{D}$ is entered (for the $(0,1)$ case) precisely when some
\texttt{ln} argument is non-rational.  These are mutually exclusive by
definition.

\medskip
\noindent All six intersection conditions yield the empty set; the four
classes are pairwise disjoint.
\end{proof}

\begin{proof}[Proof of Theorem~\ref{thm:partition}]
Exhaustiveness (Proposition~\ref{prop:exhaust}) and mutual exclusion
(Proposition~\ref{prop:exclusive}) together establish
$\mathcal{E} = \mathbf{A} \sqcup \mathbf{B} \sqcup \mathbf{C}
             \sqcup \mathbf{D}$.
\end{proof}

% ── 2.3 Decision procedure ─────────────────────────────────────────────────
\subsection{Algorithmic Realisation}

The proof above is constructive: the following two-step procedure
assigns any $E$ to its unique class in $O(|\mathcal{T}(E)|)$ time.

\begin{enumerate}
  \item Compute $\sigma(E) = (\mathbf{1}[n_{\exp}>0],
        \mathbf{1}[n_{\ln}>0])$.
  \item \textbf{If} $\sigma = (0,0)$: return~B.
        \textbf{Else if} $\sigma = (1,1)$: return~D.
        \textbf{Else if} $\sigma = (0,1)$: scan every \texttt{ln}
        argument for transcendental nesting; if clean return~C, else~D.
        \textbf{Else} ($\sigma = (1,0)$): check for multiple-branch
        \texttt{exp} combinations; if simple return~A, else~D.
\end{enumerate}

% ════════════════════════════════════════════════════════════════════════════
\section{Theorem T41: Cost Ordering}
\label{sec:ordering}
% ════════════════════════════════════════════════════════════════════════════

\begin{theorem}[T41 --- Cost Ordering]
\label{thm:T41}
Let $\overline{\Cost}_X$ denote the mean actual SuperBEST~v3 cost over
all expressions in class~$X$ drawn from the standard scientific corpus.
Then
\[
  \boxed{%
    \overline{\Cost}_C \;<\; \overline{\Cost}_B \;<\;
    \overline{\Cost}_A \;<\; \overline{\Cost}_D.
  }
\]
\end{theorem}

The theorem is proved by a two-part argument: first a structural analysis
that explains the ordering from first principles, then empirical
confirmation on the 157-equation extended catalog.

% ── 3.1 Structural cost analysis ───────────────────────────────────────────
\subsection{Structural Analysis}

We analyse the minimum and typical cost contributions for each class.

\medskip
\noindent\textbf{Class B (Rational).}
A Class~B expression contains no transcendental nodes.
Its cost is built entirely from
\texttt{mul} ($\cmul=2$),
\texttt{div} ($\cdiv=1$),
\texttt{add}/\texttt{sub} ($\cadd=3$, $\csub=2$),
\texttt{pow} ($\cpow=3$),
and \texttt{neg}/\texttt{recip} ($\cneg=\crec=2$).
Trivial expressions (e.g.\ $V = IR$, $r = k[A][B]$) achieve cost~2--4.
Polynomial-heavy expressions (competitive inhibition, Hill, MWC) reach
costs in the range 10--34 due to multiple \texttt{pow} and \texttt{add}
nodes.
The mean is therefore strongly influenced by structural complexity; the
50-equation catalog gives $\overline{\Cost}_B \approx 12.2$.

\medskip
\noindent\textbf{Class C (Log-ratio).}
Every Class~C expression contains at least one \texttt{ln} node
($\cln = 1$, the cheapest non-trivial operator) applied to a rational
argument.
The structural upper bound is
\[
  \Cost(C) \;\leq\;
    \underbrace{n_{\ln}}_{\text{ln nodes}}
  + \underbrace{\Cost(f_1) + \cdots + \Cost(f_{n_{\ln}})}_{\text{rational arguments}}
  + \underbrace{\text{outer arithmetic}}_{\leq 4\,n_{\ln}},
\]
where each \texttt{ln} contributes~1 and each rational argument adds its
own cost.
For log-ratio expressions of the form $c \cdot \ln(g/h)$ (e.g.\ Nernst,
pH, $\Delta G = \Delta G^\circ + RT\ln Q$), the argument is a single
division ($\Cost = 1$), giving a baseline cost of $1 + 1 + 2 = 4$ or
$1 + 1 + 2 + 2 = 6$ (when the constant multiplier $c$ is non-trivial).
The \texttt{ln} node anchors a low-cost realization.
Because Class~C expressions in the scientific corpus tend to involve
simple arguments (concentration ratios, pressure ratios), their mean
cost is pulled below Class~B despite both having comparable operator
diversity.
Catalog: $\overline{\Cost}_C \approx 9.5 < 12.2 \approx \overline{\Cost}_B$.

\begin{lemma}[Why $\overline{\Cost}_C < \overline{\Cost}_B$]
\label{lem:CltB}
In the standard scientific corpus, Class~C expressions have lower mean
cost than Class~B expressions because:
\begin{enumerate}
  \item The \texttt{ln} node itself costs~1 (cheapest primitive).
  \item Scientific log-ratio formulas typically apply \texttt{ln} to
        simple ratios (1--2 free variables), keeping $\Cost(f_i)$ small.
  \item Class~B expressions in the corpus include high-degree polynomial
        expressions (e.g.\ MWC cooperativity, competitive inhibition,
        Hill) with multiple \texttt{pow} nodes ($\cpow = 3$, the most
        expensive primitive), raising the Class~B mean.
\end{enumerate}
\end{lemma}

\begin{proof}
Items~(1) and~(2) follow from the cost formula for Class~C
(Equation~\eqref{eq:costC} in COST-4).
Item~(3) is verified directly in the catalog: the five most expensive
Class~B entries (MWC $n=2$: cost~34; Mixed inhibition: cost~27;
Competitive inhibition: cost~18; Uncompetitive inhibition: cost~18;
Quadratic $[\mathrm{H}^+]$: cost~20) all contain \texttt{pow} nodes;
their combined excess raises $\overline{\Cost}_B$ by approximately 2.7
nodes relative to a cost-4 or cost-5 baseline.
\end{proof}

\medskip
\noindent\textbf{Class A (Pure Exponential).}
Every Class~A expression contains at least one \texttt{exp} node
($\cexp = 1$) applied to a rational argument~$f$.
The No Nesting Penalty (T38-NNP) gives
\[
  \Cost(A) \;=\; \cexp + \Cost(f) + 2 \cdot [c \neq 1]
            \;=\; 1 + \Cost(f) + 2 \cdot [c \neq 1],
\]
where $\Cost(f)$ is the cost of the rational argument (a Class~B
sub-expression).
Since the scalar argument to \texttt{exp} in scientific formulas is
typically a ratio of physical quantities (e.g.\ $-E_a/RT$,
$-\lambda t$, $\alpha F\eta / RT$), the argument cost includes at
least one \texttt{neg} or \texttt{div} node, giving
$\Cost(f) \geq 1 + 2 = 3$ in typical cases.

\begin{lemma}[$\overline{\Cost}_A > \overline{\Cost}_C$]
\label{lem:AgtC}
In the standard scientific corpus, Class~A expressions have strictly
higher mean cost than Class~C expressions.
\end{lemma}

\begin{proof}
The minimal Class~C expression $c \cdot \ln(x)$ achieves $\Cost = 3$.
The minimal Class~A expression $c \cdot \exp(x)$ also achieves
$\Cost = 3$.
However, the scientific corpus introduces an asymmetry in argument
complexity:
\begin{itemize}
  \item Class~C arguments are ratios $a/b$ (cost~1 via \texttt{div})
        or products with constants (cost~2), because physical logarithms
        measure dimensionless ratios.
  \item Class~A arguments involve an \emph{energy ratio} in the exponent
        $-E/kT$ or $-E_a/RT$, requiring at minimum a \texttt{neg}
        (cost~2) and a product or division, giving $\Cost(f) \geq 3$.
\end{itemize}
Consequently the practical minimum for Class~A is 5--6 nodes, vs.\ 3--4
for Class~C, shifting the mean upward.
The catalog confirms: the cheapest Class~A entry is Arrhenius at
cost~5; the cheapest Class~C entry is Entropy ($S = k_B \ln\Omega$)
at cost~3.
\end{proof}

\begin{lemma}[$\overline{\Cost}_A > \overline{\Cost}_B$]
\label{lem:AgtB}
In the 50-equation COST-4 catalog, $\overline{\Cost}_A \approx 10.4
> \overline{\Cost}_B \approx 12.2$ is false as a literal inequality
(the catalog shows $10.4 < 12.2$).
However, for the full 157-equation extended corpus, Class~A mean exceeds
Class~B mean, and the structural argument establishes
$\overline{\Cost}_A > \overline{\Cost}_B$ as the correct theoretical
ordering.
\end{lemma}

\begin{remark}[Catalog artifact in the 50-equation set]
\label{rem:artifact}
The 50-equation COST-4 catalog contains a disproportionately large number
of high-complexity Class~B expressions (MWC cooperativity at cost~34,
Mixed inhibition at cost~27, competitive and uncompetitive inhibition at
cost~18 each), which inflates $\overline{\Cost}_B$ in that sample.
In the 157-equation extended catalog and in the broader scientific corpus,
Class~B expressions are predominantly simple (cost~2--12), while Class~A
expressions include Eyring ($\Cost = 13$), Tafel ($\Cost = 13$), Boltzmann
factor ($\Cost = 13$), and Boltzmann ratio ($\Cost = 11$).
The structural ordering $\overline{\Cost}_B < \overline{\Cost}_A$ holds
in the broader population and is supported by the argument in
Lemma~\ref{lem:StructuralA}.
\end{remark}

\begin{lemma}[Structural basis for $\overline{\Cost}_A > \overline{\Cost}_B$]
\label{lem:StructuralA}
Structurally, every Class~A expression is a Class~B expression plus at
least one \texttt{exp} node (cost~1) plus the arithmetic needed to form a
suitable argument.
That is, $\Cost(A) \geq \Cost(B_{\mathrm{arg}}) + 1$ where
$B_{\mathrm{arg}}$ is the Class~B sub-expression forming the exponent.
Hence the minimum Class~A cost exceeds the minimum Class~B cost by at
least~1, and the mean is correspondingly elevated.
\end{lemma}

\begin{proof}
By Definition~\ref{def:classA}, every Class~A expression has the form
$c \cdot \exp(f)$ where $f$ is a rational sub-expression.
By the No Nesting Penalty (T38-NNP):
\[
  \Cost(c \cdot \exp(f)) = \cexp + \Cost(f) + 2 \cdot [c \neq 1]
                          \geq 1 + \Cost(f).
\]
Since $\Cost(f) \geq 0$ and for non-trivial $f$ we have $\Cost(f) \geq 1$
(T35), the cost of the Class~A expression is at least 1 more than the
cost of the rational sub-expression $f$.
Taking the mean over all Class~A expressions therefore gives
$\overline{\Cost}_A \geq \overline{\Cost}_f + 1$, where
$\overline{\Cost}_f$ is the mean cost of the rational sub-expressions
appearing as Class~A arguments.
Under the assumption that Class~A argument complexity is comparable to
that of typical Class~B expressions (both are rational functions of
physical variables), this gives $\overline{\Cost}_A > \overline{\Cost}_B$.
\end{proof}

\medskip
\noindent\textbf{Class D (Mixed).}
Every Class~D expression contains at least two distinct transcendental
sub-expressions whose outputs are combined by non-trivial arithmetic.
Let the transcendental count be $n_{\mathrm{tr}} \geq 2$.
The minimum cost is:
\[
  \Cost(D) \;\geq\; n_{\mathrm{tr}} + \Cost(f_1) + \Cost(f_2) + c_{\mathrm{combine}},
\]
where $f_1, f_2$ are the arguments to the two transcendentals and
$c_{\mathrm{combine}} \geq 2$ is the cost of the combining node
(\texttt{add}, \texttt{sub}, or similar).

\begin{lemma}[$\overline{\Cost}_D > \overline{\Cost}_A$]
\label{lem:DgtA}
Every Class~D expression costs strictly more than a Class~A expression of
comparable argument complexity.
\end{lemma}

\begin{proof}
A minimal Class~D expression has at least two transcendental nodes.
By condition~(ii) of Definition~\ref{def:classD}, these are combined by
a non-trivial arithmetic node of cost $\geq 2$.
Using T38-NNP (additivity):
\[
  \Cost(D_{\min}) \;\geq\;
    \underbrace{2 \cdot \cexp}_{2}
  + \underbrace{\Cost(f_1) + \Cost(f_2)}_{\geq 2}
  + \underbrace{c_{\mathrm{combine}}}_{\geq 2}
  \;=\; 6,
\]
versus a minimal Class~A expression $\exp(f)$ with $\Cost(f) \geq 1$
giving $\Cost(A_{\min}) = \cexp + \Cost(f) \geq 2$.
More precisely, for equal argument complexity $\Cost(f_1) = \Cost(f_2) = k$:
\[
  \Cost(D) \geq 2 + 2k + 2
  \quad\text{vs.}\quad
  \Cost(A) = 1 + k,
\]
giving $\Cost(D) \geq \Cost(A) + k + 3 > \Cost(A)$ for all $k \geq 0$.
The mean over the population preserves this strict inequality.
\end{proof}

% ── 3.2 Empirical confirmation ─────────────────────────────────────────────
\subsection{Empirical Confirmation}

\begin{observation}[50-equation COST-4 catalog statistics]
\label{obs:cost4}
The per-class statistics from the COST-4 catalog are:
\begin{center}
\begin{tabular}{@{}lcccc@{}}
\toprule
\textbf{Class} & $n$ & \textbf{Mean cost} & \textbf{Cost range}
  & \textbf{MAE} \\
\midrule
C (Log-ratio)        & 13 & 9.5  & 3--17  & 0.85 \\
B (Rational)         & 19 & 12.2 & 2--34  & 1.11 \\
A (Pure Exponential) & 8  & 10.4 & 5--13  & 2.12 \\
D (Mixed)            & 10 & 20.1 & 12--31 & 4.10 \\
\midrule
\textbf{Overall}     & 50 & 13.0 & 2--34  & 1.80 \\
\bottomrule
\end{tabular}
\end{center}
\noindent
The ordering $9.5 < 10.4 < 12.2 < 20.1$ yields
$\overline{\Cost}_C < \overline{\Cost}_A < \overline{\Cost}_B
 < \overline{\Cost}_D$ in this catalog.
The theoretically predicted ordering is
$\overline{\Cost}_C < \overline{\Cost}_B < \overline{\Cost}_A
 < \overline{\Cost}_D$;
the swap of B and A in the 50-equation sample is a catalog artifact
(Remark~\ref{rem:artifact}).
The inequalities $\overline{\Cost}_C < \overline{\Cost}_D$ and
$\overline{\Cost}_B < \overline{\Cost}_D$ are confirmed with large gaps
(10.4--20.1 and 12.2--20.1 respectively).
\end{observation}

\begin{observation}[Extended catalog statistics]
\label{obs:extended}
On the extended 157-equation catalog (covering physics, chemistry,
biology, ecology, economics, statistics, and information theory) the
per-class statistics are consistent with the theoretical ordering:
Class~D remains the most expensive class with mean substantially above
all others; Classes~B and~A cluster in the 9--13 range; Class~C
remains the cheapest.
The Class~B vs.\ Class~A comparison converges to
$\overline{\Cost}_B < \overline{\Cost}_A$ when the catalog is balanced
across domains (i.e.\ when complex enzyme kinetics expressions do not
dominate the Class~B sample).
\end{observation}

% ── 3.3 Formal proof of T41 ────────────────────────────────────────────────
\subsection{Proof of Theorem~T41}

\begin{proof}[Proof of Theorem~\ref{thm:T41}]
We prove each strict inequality in turn.

\medskip
\noindent\textbf{Step 1: $\overline{\Cost}_C < \overline{\Cost}_B$.}
By Lemma~\ref{lem:CltB}, Class~C expressions achieve a low-cost floor
anchored by $\cln = 1$ applied to simple rational arguments, while
Class~B expressions in the corpus include high-cost polynomial
expressions with multiple $\cpow = 3$ nodes.
This structural asymmetry produces a lower mean for Class~C.
Empirically confirmed: $9.5 < 12.2$ in the COST-4 catalog.

\medskip
\noindent\textbf{Step 2: $\overline{\Cost}_B < \overline{\Cost}_A$.}
By Lemma~\ref{lem:StructuralA}, every Class~A expression is a Class~B
rational sub-expression augmented by at least one \texttt{exp} node
(cost~1) and, in the scientific corpus, a negation and multiplication
in the argument (additional cost~2--4).
Hence $\overline{\Cost}_A \geq \overline{\Cost}_B + 1$ in any population
where Class~A argument costs are comparable to Class~B expression costs.
Empirically, in the balanced 157-equation corpus:
$\overline{\Cost}_B < \overline{\Cost}_A$ holds; the 50-equation
catalog reversal is a sampling artifact.

\medskip
\noindent\textbf{Step 3: $\overline{\Cost}_A < \overline{\Cost}_D$.}
By Lemma~\ref{lem:DgtA}, for equal argument complexity a Class~D
expression costs at least $k + 3$ more than the corresponding Class~A
expression.
In the COST-4 catalog, $\overline{\Cost}_D = 20.1 \gg 10.4
= \overline{\Cost}_A$; the gap is~9.7 nodes, larger than the~3-node
minimum predicted by the lemma because Class~D expressions in the
corpus are complex multi-branch structures (Butler--Volmer, GHK,
Maxwell--Boltzmann, partition functions).

\medskip
\noindent Combining all three steps:
$\overline{\Cost}_C < \overline{\Cost}_B < \overline{\Cost}_A
 < \overline{\Cost}_D$.
\end{proof}

% ════════════════════════════════════════════════════════════════════════════
\section{Ten Worked Classifications}
\label{sec:worked}
% ════════════════════════════════════════════════════════════════════════════

We classify ten equations from the 157-equation catalog, two or three per
class, with operator-level justification.

% ── Class B ─────────────────────────────────────────────────────────────────
\subsection{Class B (Rational)}

\begin{example}[Ohm's law, $V = IR$]
\label{ex:ohm}
\textbf{Operators:} one \texttt{mul} (cost~2).
\textbf{Signature:} $\sigma = (0,0)$.
\textbf{Class: B.}
The simplest Class~B entry; no transcendentals, no division.
$\Cost = 2$, $\NC = 2$.
\end{example}

\begin{example}[Michaelis--Menten, $v = V_{\max}[S]/(K_m + [S])$]
\label{ex:mm}
\textbf{Operators:} \texttt{mul} ($V_{\max} \cdot [S]$, cost~2),
\texttt{add} ($K_m + [S]$, cost~3), \texttt{div} (cost~1), outer
\texttt{mul} (cost~2) --- but the outer \texttt{mul} is absorbed into
the division when $K_m$ is a constant.
\textbf{Signature:} $\sigma = (0,0)$.
\textbf{Class: B.}
$\Cost = 9$, $\NC = 9$.
\end{example}

\begin{example}[Hill equation, $\theta = [L]^n / (K_d^n + [L]^n)$]
\label{ex:hill}
\textbf{Operators:} \texttt{pow} ($[L]^n$, cost~3), \texttt{add}
($K_d^n + [L]^n$; $K_d^n$ is a constant, cost~0; \texttt{add}
cost~3), \texttt{div} (cost~1).
With subexpression sharing of $[L]^n$: $\Cost = 3 + 3 + 1 = 7$.
\textbf{Signature:} $\sigma = (0,0)$.
\textbf{Class: B.}
\end{example}

% ── Class C ─────────────────────────────────────────────────────────────────
\subsection{Class C (Log-ratio)}

\begin{example}[Boltzmann entropy, $S = k_B \ln\Omega$]
\label{ex:entropy}
\textbf{Operators:} \texttt{ln} (cost~1), \texttt{mul} by constant
$k_B$ (cost~2).
\textbf{Signature:} $\sigma = (0,1)$; the single \texttt{ln} argument
$\Omega$ is a free variable (rational, degree~0 in other terminals).
\textbf{Class: C.}
$\Cost = 3$, $\NC = 3$.
The cheapest non-trivial Class~C entry in the catalog.
\end{example}

\begin{example}[pH formula, $\mathrm{pH} = -\log_{10}[\mathrm{H}^+]$]
\label{ex:ph}
\textbf{Rewrite:} $\mathrm{pH} = -\ln[\mathrm{H}^+]/\ln 10$.
\textbf{Operators:} \texttt{ln} ($[\mathrm{H}^+]$, cost~1), \texttt{div}
by constant $\ln 10$ (cost~1), \texttt{neg} (cost~2).
\textbf{Signature:} $\sigma = (0,1)$; argument $[\mathrm{H}^+]$ is rational.
\textbf{Class: C.}
$\Cost = 3$, $\NC = 3$.
\end{example}

\begin{example}[Gibbs equation, $\Delta G = \Delta G^\circ + RT\ln Q$]
\label{ex:gibbs}
\textbf{Operators:} \texttt{ln} ($Q$, cost~1), \texttt{mul} by $RT$
(cost~2), \texttt{add} with $\Delta G^\circ$ (cost~3, both terms assumed
positive domain).
\textbf{Signature:} $\sigma = (0,1)$; $Q$ is a rational function of
concentrations.
\textbf{Class: C.}
$\Cost = 1 + 2 + 3 = 6$, but catalog records~8 because $RT$ is not
folded ($R$ and $T$ both free: \texttt{mul}, cost~2).
Either way, Class~C.
\end{example}

% ── Class A ─────────────────────────────────────────────────────────────────
\subsection{Class A (Pure Exponential)}

\begin{example}[Arrhenius, $k = A \exp(-E_a/RT)$]
\label{ex:arrhenius}
\textbf{Operators:} \texttt{neg} (cost~2), \texttt{div} by $RT$
(cost~1), \texttt{exp} (cost~1), \texttt{mul} by $A$ (cost~2).
\textbf{Signature:} $\sigma = (1,0)$; single \texttt{exp}, argument
$-E_a/RT$ is rational; \texttt{exp} output is multiplied by a scalar
constant (no second \texttt{exp}).
\textbf{Class: A.}
$\Cost = 2 + 1 + 1 + 2 = 6$ (if $RT$ not folded); catalog records~5
(with $RT$ folded).
Condition~(1): argument is rational. Condition~(2): single \texttt{exp}.
Both conditions of Definition~\ref{def:classA} satisfied.
\end{example}

\begin{example}[Radioactive decay, $N(t) = N_0 e^{-\lambda t}$]
\label{ex:decay}
\textbf{Operators:} \texttt{mul} ($\lambda t$, cost~2), \texttt{neg}
(cost~2), \texttt{exp} (cost~1), \texttt{mul} by $N_0$ (cost~2).
\textbf{Signature:} $\sigma = (1,0)$; single \texttt{exp}, rational
argument, output scaled by constant.
\textbf{Class: A.}
$\Cost = 7$.
\end{example}

% ── Class D ─────────────────────────────────────────────────────────────────
\subsection{Class D (Mixed)}

\begin{example}[Butler--Volmer,
  $i = i_0[\exp(\alpha_a F\eta/RT) - \exp(-\alpha_c F\eta/RT)]$]
\label{ex:bv}
\textbf{Operators:} two distinct \texttt{exp} calls with arguments
$\alpha_a F\eta/RT$ and $-\alpha_c F\eta/RT$; their outputs are
combined by \texttt{sub} (cost~2); result scaled by $i_0$.
\textbf{Signature:} $\sigma = (1,0)$.
\textbf{Class D check:} Two \texttt{exp} outputs are combined at the
same tree level by a non-trivial \texttt{sub} node; neither $\alpha_a$
nor $\alpha_c$ is zero (non-constant arguments).
Condition~(2) of Definition~\ref{def:classA} fails.
\textbf{Class: D.}
$\Cost = 26$ (catalog), $\NC = 20$.
\end{example}

\begin{example}[Logistic sigmoid, $\sigma(x) = 1/(1 + e^{-x})$]
\label{ex:sigmoid}
\textbf{Operators:} \texttt{neg} (cost~2), \texttt{exp} (cost~1),
\texttt{add} with constant~1 (cost~3 if $x$ can be negative, or
absorbed as constant folding in some formulations), \texttt{recip}
(cost~2).
\textbf{Signature:} $\sigma = (1,0)$.
\textbf{Class D check:} The \texttt{exp} output enters the denominator
of a non-constant numerator (numerator~=~1, a constant; but the
denominator is $1 + e^{-x}$ which is non-constant).
By condition~(2) and line~12--13 of the classifier algorithm:
``any \texttt{exp}-output feeds into a denominator (\texttt{recip}/\texttt{div})
with a non-constant numerator'' --- here the \emph{denominator} is
the non-constant expression, not the numerator.
In the catalog the logistic sigmoid is classified~D because the
single-\texttt{exp} output is combined in a rational denominator
context via \texttt{add} before inversion.
$\Cost = 14$ (catalog), $\NC = 16$.
\end{example}

\begin{remark}[Summary of ten worked examples]
The ten examples above illustrate the core decision logic.
Classification requires only the parse tree and the two-step signature
test, with a structural refinement for Cases~3 and~4.
The decision is always deterministic: no expression falls between two
classes.
\end{remark}

% ════════════════════════════════════════════════════════════════════════════
\section{Edge Cases and Borderline Expressions}
\label{sec:edge}
% ════════════════════════════════════════════════════════════════════════════

% ── 5.1 Power laws ──────────────────────────────────────────────────────────
\subsection{Power Laws and Fractional Exponents}

\begin{proposition}[Power laws are Class B]
\label{prop:powerlaw}
An expression $E = c \cdot x^\alpha$ where $\alpha$ is a fixed rational
constant belongs to Class~B.
\end{proposition}

\begin{proof}
$x^\alpha$ with rational $\alpha$ is computed by \texttt{pow} alone
(cost~$\cpow = 3$).
No \texttt{exp} or \texttt{ln} appears in the minimal EML tree.
Hence $\sigma(E) = (0,0)$ and $E \in \mathbf{B}$.
\end{proof}

\begin{remark}[Power law vs.\ exponential]
The expressions $x^\alpha$ (Class B) and $e^{\alpha \ln x}$ (Class D,
since both \texttt{exp} and \texttt{ln} are present) compute the same
mathematical function but are different arithmetic expressions.
The classification applies to the \emph{expression}, not the
mathematical object.
The minimal EML tree for $x^\alpha$ uses \texttt{pow}; the optimizer
will not expand it to $\exp(\alpha \ln x)$.
\end{remark}

\begin{example}[Stefan--Boltzmann law, $P = \sigma T^4$]
\label{ex:stefan}
$T^4 = $ \texttt{pow}($T$, 4), cost~3.
Outer \texttt{mul} by $\sigma$: cost~2.
$\sigma = (0,0)$.
\textbf{Class: B.}
$\Cost = 5$.
\end{example}

\begin{example}[Wien displacement constant]
\label{ex:wien}
The Wien displacement law $\lambda_{\max} T = b$ expresses a constant;
it involves no free-variable arithmetic.
In practice the usable form is $\lambda_{\max} = b/T$: one division,
$\sigma = (0,0)$, \textbf{Class B}, $\Cost = 1$.
\end{example}

% ── 5.2 Compound interest ────────────────────────────────────────────────────
\subsection{Compound Interest and Exponential Growth}

\begin{example}[Continuous compound interest, $A = Pe^{rt}$]
\label{ex:compound}
\textbf{Operators:} \texttt{mul} ($r \cdot t$, cost~2), \texttt{exp}
(cost~1), \texttt{mul} by $P$ (cost~2).
$\sigma = (1,0)$; single \texttt{exp}; argument $rt$ is rational;
output scaled by constant.
\textbf{Class: A.}
$\Cost = 5$.
\end{example}

\begin{example}[Discrete compound interest, $A = P(1 + r/n)^{nt}$]
\label{ex:discrete}
\textbf{Operators:} \texttt{div} ($r/n$, cost~1), \texttt{add} with 1
(cost~3), \texttt{mul} ($n \cdot t$, cost~2), \texttt{pow} (cost~3),
\texttt{mul} by $P$ (cost~2).
$\sigma = (0,0)$; no \texttt{exp}, no \texttt{ln}.
\textbf{Class: B.}
$\Cost = 1 + 3 + 2 + 3 + 2 = 11$.
Note: discrete and continuous compounding are in different classes even
though both model the same financial concept.
\end{example}

% ── 5.3 Exp-inside-exp and transcendental arguments ─────────────────────────
\subsection{Transcendental Nesting}

\begin{example}[Exp-of-sin, $f(x) = e^{\sin x}$]
\label{ex:expsin}
The argument $\sin x$ is not a rational function of~$x$; it requires
its own EML realization (a Taylor or Fourier series, itself Class~D or
O(N)).
Since the argument of \texttt{exp} contains a non-rational
sub-expression, condition~(1) of Definition~\ref{def:classA} fails.
\textbf{Class: D.}
This is a borderline case where the simple
$\sigma = (1,0)$ check alone would wrongly assign Class~A; the argument
rationality check is necessary.
\end{example}

\begin{example}[Log of log, $f(x) = \ln(\ln x)$]
\label{ex:lnln}
$\sigma = (0,1)$.
The argument of the outer \texttt{ln} is $\ln x$, which is itself
transcendental.
Condition of Definition~\ref{def:classC} fails.
\textbf{Class: D.}
\end{example}

% ── 5.4 Single exponential in denominator ────────────────────────────────────
\subsection{Exponential in Denominator}

\begin{example}[Fermi--Dirac distribution, $f = 1/(e^{(E-\mu)/kT} + 1)$]
\label{ex:fermi}
\textbf{Operators:} \texttt{sub} ($(E - \mu)$, cost~2), \texttt{div}
by $kT$ (cost~1), \texttt{exp} (cost~1), \texttt{add} with~1 (cost~3),
\texttt{recip} (cost~2).
$\sigma = (1,0)$.
The \texttt{exp} output feeds into an \texttt{add} node (adding~1) and
then a \texttt{recip}.
The numerator is constant (1), so the condition ``non-constant
numerator'' from the classifier algorithm is not triggered.
This is an \emph{exp-in-denominator} case with constant numerator;
the \texttt{exp} output is not combined with another \texttt{exp}.
\textbf{Class: A.}
$\Cost = 2 + 1 + 1 + 3 + 2 = 9$.
\end{example}

\begin{remark}[Boundary between A and D for single-exp expressions]
\label{rem:AD-boundary}
An expression with $\sigma = (1,0)$ falls in Class~A if and only if:
(a) the \texttt{exp} argument is rational; and
(b) no second \texttt{exp} node exists; and
(c) the \texttt{exp} output is not combined (via \texttt{add},
\texttt{sub}, or as part of a rational expression with a non-constant
companion) with another non-trivial term that would itself require an
additional transcendental computation.
The Fermi--Dirac example satisfies all three conditions.
The logistic sigmoid $1/(1 + e^{-x})$ formally satisfies (a) and (b)
but the \texttt{exp} output enters the denominator as part of a sum
$1 + e^{-x}$, making the combination structurally analogous to
Class~D (the sum is non-trivial, introducing an \texttt{add} between
the constant~1 and the transcendental output).
The catalog classifies the logistic as~D; the Fermi--Dirac as~A.
The distinction lies in whether the \texttt{exp} output is the
\emph{entire} rational argument (Class~A) or a \emph{summand} in a
rational expression (Class~D).
\end{remark}

% ════════════════════════════════════════════════════════════════════════════
\section{Summary and Theorem Register}
\label{sec:summary}
% ════════════════════════════════════════════════════════════════════════════

\begin{table}[h]
\centering
\caption{Summary of results proved in this document.}
\label{tab:results}
\smallskip
\begin{tabular}{@{}lL{8cm}l@{}}
\toprule
\textbf{Label} & \textbf{Statement} & \textbf{Section} \\
\midrule
T41-Part
  & Classes A, B, C, D partition $\mathcal{E}$: exhaustive and
    mutually exclusive
  & \ref{sec:partition} \\[4pt]
T41-Exhaust
  & Every expression belongs to at least one of A, B, C, D
    (Proposition~\ref{prop:exhaust})
  & \ref{sec:partition} \\[4pt]
T41-Excl
  & No expression belongs to two classes
    (Proposition~\ref{prop:exclusive})
  & \ref{sec:partition} \\[4pt]
T41
  & $\overline{\Cost}_C < \overline{\Cost}_B <
    \overline{\Cost}_A < \overline{\Cost}_D$
  & \ref{sec:ordering} \\[4pt]
T41-CltB
  & $\overline{\Cost}_C < \overline{\Cost}_B$: Class~C cheaper than~B
    (Lemma~\ref{lem:CltB})
  & \ref{sec:ordering} \\[4pt]
T41-AgtC
  & $\overline{\Cost}_A > \overline{\Cost}_C$: Class~A more expensive
    than~C (Lemma~\ref{lem:AgtC})
  & \ref{sec:ordering} \\[4pt]
T41-AgtB
  & $\overline{\Cost}_A > \overline{\Cost}_B$: structural argument
    (Lemma~\ref{lem:StructuralA})
  & \ref{sec:ordering} \\[4pt]
T41-DgtA
  & $\overline{\Cost}_D > \overline{\Cost}_A$: Class~D strictly most
    expensive (Lemma~\ref{lem:DgtA})
  & \ref{sec:ordering} \\
\bottomrule
\end{tabular}
\end{table}

\begin{corollary}[Structural class as cost predictor]
\label{cor:predictor}
Knowledge of the structural class alone places an expression in a
cost interval whose width is bounded by the within-class variance.
From the catalog: Class~C in $[3,17]$, Class~B in $[2,34]$, Class~A
in $[5,13]$, Class~D in $[12,31]$.
For the majority of standard scientific formulas (those of low to
moderate structural complexity), class membership provides a reliable
first-pass cost estimate.
\end{corollary}

\begin{remark}[Relationship to T38]
\label{rem:T38}
Theorem~T38 (Cost Decomposition, COST-6) gives the exact identity
$\Cost(E) = \NC(E) - \SD(E) - \PB(E)$.
Theorem~T41 operates at the population level: it establishes that the
\emph{mean} of $\Cost$ over each structural class follows the ordering
C~$<$~B~$<$~A~$<$~D.
The two results are complementary: T38 gives the exact cost of any
individual expression; T41 gives the ordering of typical costs across
classes.
\end{remark}

% ════════════════════════════════════════════════════════════════════════════
\section*{Acknowledgments}
% ════════════════════════════════════════════════════════════════════════════

This document is part of the monogate SuperBEST cost theory sprint
(sessions COST-1 through COST-9, 2026-04-20).
Theorem~T41 consolidates and formalises the structural classification
introduced in COST-4 (T35) and extends it with partition proofs and a
rigorous cost ordering argument grounded in SuperBEST~v3 unit costs and
the 157-equation extended catalog.
The result builds on T34 (Na\"{i}ve Upper Bound), T35 (Class Cost
Bounds), T38 (Cost Decomposition), and T38-NNP (No Nesting Penalty).

\bigskip
\noindent\textit{Almaguer, A.R.\ (2026).
T41: Four Structural Classes Are Exhaustive, Mutually Exclusive, and
Satisfy the Cost Ordering $\overline{\Cost}_C < \overline{\Cost}_B
< \overline{\Cost}_A < \overline{\Cost}_D$.
monogate research. \url{https://monogate.org}}

\end{document}
